\clearpage
\VerdierPhysicalPageAnchor{0078}%
\VerdierSemanticAnchor{ch1-2-5-4}
\noindent\textbf{2.5.4.}\quad
运算 $\star$ 与“同伦之和”具有下列性质：

\VerdierSemanticAnchor{ch1-2-5-4-eq1}
\begin{equation}
  \begin{aligned}
    u\star(b+a) &= u\star b+u\star a,\\
    (b+a)\star v &= b\star v+a\star v,\\
    uu'\star a &= u\star(u'\star a),\\
    a\star vv' &= (a\star v)\star v'.
  \end{aligned}
  \tag{2.5.4.1}
\end{equation}

\VerdierSemanticAnchor{ch1-2-5-5}
\noindent\textbf{2.5.5.}\quad
若存在一个连接两个态射 $f,g\colon X\rightrightarrows Y$ 的同伦，
则称 $f$ 与 $g$ \emph{同伦}。

\VerdierSemanticAnchor{ch1-2-5-5-difference-criterion}
若态射 $f-g$ 同伦于零，则态射 $f,g$ 同伦。

\VerdierSemanticAnchor{ch1-2-5-5-shift-stability}
设 $f\colon X\longrightarrow Y$ 是一个同伦于零的态射。对任意
$\sigma\in\mathbb Z[n]$，态射 $T(\sigma)f$ 同伦于零。

\VerdierSemanticAnchor{ch1-2-5-5-null-homotopic-ideal}
所有同伦于零的态射构成 $\operatorname{comp}^{n}(C)$ 的一个理想，
即一个同伦于零的态射与任意态射无论按哪种次序复合，所得态射都同伦于零，
而两个同伦于零的态射之直和也同伦于零。

\VerdierSemanticAnchor{ch1-2-5-6}
\noindent\textbf{2.5.6.}\quad
设 $\varphi\colon[n]\longrightarrow[m]$ 是一个映射。

\VerdierSemanticAnchor{ch1-2-5-6-integration-preserves-homotopy}
函子 $\int_\varphi^\Sigma$ 与 $\int_\varphi^\Pi$ 将
$\operatorname{comp}^{n}(C)$ 中的每对同伦态射送到
$\operatorname{comp}^{m}(C)$ 中的一对同伦态射。

\VerdierSemanticAnchor{ch1-2-5-6-multiadditive-setup}
同样，设
$F\colon\prod_{i\in[n]}C_i\longrightarrow C$ 为一个多加法函子，

\VerdierSemanticAnchor{ch1-2-5-6-fn-functor}
\begin{equation*}
  F^{[n]}\colon
  \prod_{i\in[n]}\operatorname{comp}(C_i)
  \longrightarrow \operatorname{comp}^{n}(C)
\end{equation*}
为 (2.4.2.4) 中所述的函子；设 $X_i$ 是
$\operatorname{comp}(C_i)$ 的对象 $(i\ne j)$，并设
$f,g\colon X\rightrightarrows Y$ 是 $\operatorname{comp}(C_j)$
中的两个同伦态射。

\VerdierSemanticAnchor{ch1-2-5-6-fn-preserves-homotopy}
$\operatorname{comp}^{n}(C)$ 中的态射
\begin{equation*}
  F^{[n]}(X_1,\ldots,f,\ldots,X_n)
  \quad\text{与}\quad
  F^{[n]}(X_1,\ldots,g,\ldots,X_n)
\end{equation*}
彼此同伦。

\VerdierSemanticAnchor{ch1-2-5-6-totalizations-preserve-homotopy}
因此，$\operatorname{comp}(C)$ 中的态射
\begin{equation*}
  \operatorname{comp}_{\Pi}F(X_1,\ldots,f,\ldots,X_n)
  \quad\text{与}\quad
  \operatorname{comp}_{\Pi}F(X_1,\ldots,g,\ldots,X_n)
\end{equation*}
彼此同伦（对 $\operatorname{comp}_{\Sigma}F$ 也有相同结论）。

\VerdierSemanticAnchor{ch1-2-5-6-reference}
我们把所有这些断言的验证留给读者，读者也可参见 [1]。

\VerdierSemanticAnchor{ch1-2-5-7}
\noindent\textbf{2.5.7.}\quad
用 $K(C)$ 表示下列范畴：

\VerdierSemanticAnchor{ch1-2-5-7-objects}
\noindent\hspace*{2em}--\ $K(C)$ 的对象是 $\operatorname{comp}(C)$ 的对象。
